%=====================================================================

\section{The Kohler--Jobin inequality and the transfer to Faber--Krahn}
\label{sec:kj}

This section is logically self-contained. In Part~1 we reconstruct from
scratch, by the Kohler--Jobin rearrangement technique, the spectral shape
inequality
\[
  T(\Om)\,\lambda_1(\Om)^{(n+2)/2}\ \ge\ T(B)\,\lambda_1(B)^{(n+2)/2}
  \qquad(B\text{ any ball}),
\]
with equality only for balls. In Part~2 we feed this inequality, together
with the global sharp Saint--Venant stability estimate proved elsewhere in
the manuscript, into a short and fully explicit computation that yields
sharp quantitative Faber--Krahn stability,
$\delta_{\mathrm{FK}}(\Om)\ge c_{\mathrm{FK}}(n)\,\Fr(\Om)^2$.

The technique, due to Kohler--Jobin~\cite{KJ1,KJ2,KJ3}, is a spherical
rearrangement that, unlike Schwarz symmetrisation, preserves not the
\emph{volume} of the superlevel sets but a \emph{torsional} functional of
them. We follow the modern exposition of Brasco~\cite[\S\S3--5,
Thm.~1.1]{Brasco2014}, specialised throughout to the linear case (the
Laplacian, $p=2$, and the Dirichlet energy quotient).

\subsection{Notation and scaling conventions}
\label{ssec:kj-notation}

Throughout, $n\ge 2$ and $\Om\subset\R^n$ is open with finite measure. We
use:
\begin{itemize}
\item $u_\Om\in H^1_0(\Om)$, the \emph{torsion function}, the unique weak
solution of $-\Delta u_\Om=1$ in $\Om$, $u_\Om=0$ on $\partial\Om$;
\item $T(\Om)=\int_\Om u_\Om\dx$, the \emph{torsional rigidity}, and
$E(\Om)=-T(\Om)/2$;
\item $\lambda_1(\Om)$, the first Dirichlet eigenvalue,
\begin{equation}\label{eq:rayleigh}
  \lambda_1(\Om)=\min_{\,0\ne v\in H^1_0(\Om)}
  \frac{\int_\Om\abs{\nabla v}^2\dx}{\int_\Om v^2\dx},
\end{equation}
attained by a first eigenfunction $u>0$ solving $-\Delta u=\lambda_1(\Om)u$;
\item $B^*$, the ball centred at $0$ with $\abs{B^*}=\abs\Om$; $B_r$ the
open ball of radius $r$ centred at $0$; $B_1$ the unit ball;
$\omega_n=\abs{B_1}$;
\item $L_n=\lambda_1(B_1)=j_{n/2-1,1}^2$ and
$T_n=T(B_1)=\omega_n/\bigl(n(n+2)\bigr)$;
\item $\Fr(\Om)=\inf_{x_0\in\R^n}\abs{\Om\,\triangle\,(x_0+B^*)}/\abs{B^*}
\in[0,2)$, the Fraenkel asymmetry;
\item $\delta_{\mathrm{FK}}(\Om)
=(\abs\Om/\omega_n)^{2/n}\bigl(\lambda_1(\Om)-\lambda_1(B^*)\bigr)$, the
scale-invariant Faber--Krahn deficit.
\end{itemize}

The torsion function of a ball is $u_{B_r}(x)=(r^2-\abs x^2)/(2n)$, whence
\begin{equation}\label{eq:torsionball}
  T(B_r)=\frac{\omega_n}{n(n+2)}\,r^{\,n+2}=T_n\,r^{\,n+2}.
\end{equation}
Under dilation $T(t\Om)=t^{n+2}T(\Om)$ and $\lambda_1(t\Om)=t^{-2}\lambda_1(\Om)$,
so
\begin{equation}\label{eq:psi-def}
  \Psi(\Om):=T(\Om)\,\lambda_1(\Om)^{(n+2)/2}
\end{equation}
is scale invariant; in particular $\Psi(B_r)=\Psi(B_1)=T_n L_n^{(n+2)/2}$
for every ball. We abbreviate $\beta:=(n+2)/2$.

%---------------------------------------------------------------------
\subsection{Part 1: the Kohler--Jobin inequality, proved from scratch}
\label{ssec:kj-proof}

\begin{theorem}[Kohler--Jobin]\label{thm:kj}
For every open $\Om\subset\R^n$ of finite measure with $\lambda_1(\Om)<\infty$
and every ball $B\subset\R^n$,
\begin{equation}\label{eq:kj}
  T(\Om)\,\lambda_1(\Om)^{\beta}\ \ge\ T(B)\,\lambda_1(B)^{\beta},
  \qquad\beta=\tfrac{n+2}{2}.
\end{equation}
Equivalently $\Psi(\Om)\ge\Psi(B_1)=T_n L_n^{\beta}$, with equality if and
only if $\Om$ is a ball (up to translation and a Lebesgue-null modification).
\end{theorem}

\subsubsection*{Strategy}

Both sides of \eqref{eq:kj} have the same scaling and $\Psi(B)=\Psi(B_1)$
for every ball; so it suffices to compare $\Om$ with a single, well-chosen
ball. The idea is to rearrange the first eigenfunction $u$ of $\Om$ into a
radial decreasing function $u^*$ on a ball $B$ in such a way that
\begin{enumerate}[label=\textup{(\Alph*)}]
\item\label{prop:A} the Dirichlet energy is \emph{exactly preserved},
$\int_B\abs{\nabla u^*}^2=\int_\Om\abs{\nabla u}^2$, and
\item\label{prop:B} the $L^2$ norm \emph{does not decrease},
$\int_B(u^*)^2\ge\int_\Om u^2$.
\end{enumerate}
Granting \ref{prop:A}--\ref{prop:B}, the Rayleigh principle
\eqref{eq:rayleigh} on $B$ gives
\begin{equation}\label{eq:rayleigh-chain}
  \lambda_1(B)\ \le\
  \frac{\int_B\abs{\nabla u^*}^2}{\int_B(u^*)^2}
  \ \le\
  \frac{\int_\Om\abs{\nabla u}^2}{\int_\Om u^2}
  =\lambda_1(\Om),
\end{equation}
the last equality because $u$ is the first eigenfunction of $\Om$.

The decisive design choice is which geometric quantity of the superlevel
sets the ball $B$ is built to match. Schwarz symmetrisation matches
\emph{volume} and yields the Faber--Krahn inequality directly; the
Kohler--Jobin rearrangement instead matches a \emph{torsional functional},
the \emph{modified torsional rigidity} introduced next, and this is exactly
what allows the Dirichlet energy to be preserved (not merely decreased)
while the $L^2$ norm grows. The relation $\lambda_1(B)\le\lambda_1(\Om)$
\emph{at equal torsion} of $B$ and $\Om$ is then the inequality
\eqref{eq:kj}.

\subsubsection*{Standing regularity facts (cited, not proved)}
\begin{enumerate}[label=\textup{(R\arabic*)}]
\item\label{R:reg} \emph{(Regularity of the eigenfunction.)} The first
eigenfunction $u$ of $\Om$ may be taken positive, solves
$-\Delta u=\lambda u$ in $\Om$ with $\lambda=\lambda_1(\Om)$, and satisfies
$u\in L^\infty(\Om)$ (Moser iteration) and $u\in C^\infty(\Om)$ in the
interior: indeed $-\Delta u=\lambda u$ with $u\in L^\infty$ gives $u\in
C^{1,\alpha}_{\rm loc}$ by Schauder, and bootstrapping the equation gives
$u\in C^\infty(\Om)$ (\cite[Prop.~2.4]{Brasco2014}, \cite{GilbargTrudinger2001}).
In particular Sard's theorem applies to $u$ in the interior, so a.e.\
$t\in(0,M)$ is a regular value. Put $M:=\norm{u}_{L^\infty(\Om)}<\infty$.
\item\label{R:coarea} \emph{(Coarea formula.)} For $u\in W^{1,1}_{\rm loc}$
and Borel $g\ge0$,
$\int_\Om g\,\abs{\nabla u}\dx=\int_0^M\bigl(\int_{\{u=t\}}g\dH\bigr)\dt$
(Federer; see Maggi, \emph{Sets of Finite Perimeter}, Thm.~13.1).
\item\label{R:saint-venant} \emph{(Saint--Venant inequality.)} For any open
bounded $A\subset\R^n$, $T(A)\le T(A^\#)$ where $A^\#$ is the ball with
$\abs{A^\#}=\abs A$; in the scale-free form we use it,
\begin{equation}\label{eq:sv-scale}
  T(A)\ \le\ T_n\Bigl(\frac{\abs A}{\omega_n}\Bigr)^{(n+2)/n},
\end{equation}
with equality iff $A$ is a ball. (\cite[(1.7) with $p=2$]{Brasco2014}.)
\item\label{R:absc} \emph{(Distribution function.)} The distribution
function $\mu_u(t):=\abs{\{u>t\}}$ is non-increasing; together with the
coarea formula this gives, for a.e.\ regular value $t$, that
$\int_{\{u=t\}}\abs{\nabla u}^{-1}\dH$ is finite and
$\mu_u'(t)=-\int_{\{u=t\}}\abs{\nabla u}^{-1}\dH$ (Brothers--Ziemer).
\end{enumerate}

\subsubsection*{Two preliminary identities}

We record the variational form of $T$ and the level-set flux identity for
the eigenfunction; both are elementary.

\begin{lemma}[Variational torsion and the eigenfunction flux]\label{lem:prelim}
\leavevmode
\begin{enumerate}[label=\textup{(\roman*)}]
\item\label{it:Fvar} For open $A$ of finite measure, with
$F(w):=\int_A w\dx-\tfrac12\int_A\abs{\nabla w}^2\dx$,
\begin{equation}\label{eq:Tvar}
  T(A)=2\max_{w\in H^1_0(A)}F(w),
\end{equation}
the maximiser being the torsion function $u_A$; equivalently
$T(A)=\max\{(\int_A w)^2/\int_A\abs{\nabla w}^2:0\ne w\in H^1_0(A)\}$.
\item\label{it:flux} The first eigenfunction $u$ satisfies, for a.e.\
$t\in(0,M)$,
\begin{equation}\label{eq:flux}
  \int_{\{u=t\}}\abs{\nabla u}\dH=\lambda\int_{\{u>t\}}u\dx .
\end{equation}
\end{enumerate}
\end{lemma}

\begin{proof}
\ref{it:Fvar}: $F$ is strictly concave and coercive on $H^1_0(A)$; its Euler
equation is $-\Delta w=1$, so the maximiser is $u_A$ and
$F(u_A)=\int u_A-\tfrac12\int\abs{\nabla u_A}^2=\int u_A-\tfrac12\int u_A
=\tfrac12\int u_A=\tfrac12 T(A)$, using $\int\abs{\nabla u_A}^2=\int u_A$
(test the equation by $u_A$). The second form is the standard Rayleigh-type
characterisation of $T$ (\cite[Prop.~2.2]{Brasco2014} with $p=2$).

\ref{it:flux}: integrate $-\Delta u=\lambda u$ over the
(a.e.\ smooth) superlevel set $\{u>t\}$ and use the divergence theorem with
outward normal $\nu=-\nabla u/\abs{\nabla u}$ on $\partial\{u>t\}=\{u=t\}$:
$\int_{\{u=t\}}\abs{\nabla u}\dH=-\int_{\{u=t\}}\partial_\nu u\dH
=\int_{\{u>t\}}(-\Delta u)\dx=\lambda\int_{\{u>t\}}u\dx$.
\end{proof}

\subsubsection*{The modified torsional rigidity}

The engine of the technique is the following functional, which assigns to a
nonnegative function a torsion-type quantity that is sensitive to the
\emph{level structure} of the function, not just to its domain.

\begin{definition}[Modified torsional rigidity]\label{def:tmod}
Let $u\in H^1_0(\Om)\cap L^\infty(\Om)$, $u\ge0$, $M=\norm u_\infty$, with
distribution function $\mu(t)=\abs{\{u>t\}}$. Assume the integrability
condition
\begin{equation}\label{eq:refcond}
  t\ \longmapsto\ \frac{\mu(t)}{\int_{\{u=t\}}\abs{\nabla u}\dH}\ \in\ L^\infty(0,M)
\end{equation}
For the first eigenfunction this holds: by \eqref{eq:flux} the ratio in
\eqref{eq:refcond} equals $\mu(t)/(\lambda\int_{\{u>t\}}u)$. On $[M/2,M)$ the
bound $\int_{\{u>t\}}u\ge t\,\mu(t)$ gives ratio $\le 1/(\lambda t)\le
2/(\lambda M)$; on $(0,M/2]$ the monotonicity $\int_{\{u>t\}}u\ge
\int_{\{u>M/2\}}u>0$ together with $\mu(t)\le\abs\Om$ gives ratio $\le
\abs\Om/(\lambda\int_{\{u>M/2\}}u)$. Hence the ratio is bounded on $(0,M)$
(cf.~\cite[Lem.~3.2]{Brasco2014}). With
$\mathcal L:=\{g\in\mathrm{Lip}([0,M]):g(0)=0\}$, define the
\emph{modified torsional rigidity of $\Om$ along $u$}
\begin{equation}\label{eq:tmod-def}
  \Tmod(\Om;u):=2\,\sup_{g\in\mathcal L}F\bigl(g\circ u\bigr),
  \qquad F(w)=\int_\Om w\dx-\tfrac12\int_\Om\abs{\nabla w}^2\dx .
\end{equation}
\end{definition}

Since the competitors $g\circ u$ form a subclass of $H^1_0(\Om)$, comparison
with \eqref{eq:Tvar} gives $\Tmod(\Om;u)\le T(\Om)$ always.

\begin{lemma}[Closed form of $\Tmod$]\label{lem:tmod-formula}
Under the hypotheses of Definition~\ref{def:tmod}, writing
$P(t):=\int_{\{u=t\}}\abs{\nabla u}\dH$, the supremum in
\eqref{eq:tmod-def} is uniquely attained by the nondecreasing $g_0$ with
$g_0'(t)=\mu(t)/P(t)$, and
\begin{equation}\label{eq:tmod-formula}
  \Tmod(\Om;u)=\int_0^M\frac{\mu(t)^2}{P(t)}\dt .
\end{equation}
\end{lemma}

\begin{proof}
We may restrict to nondecreasing $g$: replacing $g$ by
$\tilde g(t)=\int_0^t\abs{g'}$ leaves $\int_\Om\abs{\nabla(g\circ u)}^2$
unchanged (since $\tilde g'=\abs{g'}$ a.e.) and increases $\int_\Om g\circ
u$ (since $\tilde g\ge g$), hence increases $F$. For nondecreasing
$g\in\mathcal L$, Cavalieri's principle gives
$\int_\Om g(u)\dx=\int_0^M g'(t)\mu(t)\dt$, while the coarea formula
\ref{R:coarea} gives
$\int_\Om\abs{\nabla(g\circ u)}^2\dx=\int_0^M g'(t)^2 P(t)\dt$. Thus
\[
  F(g\circ u)=\int_0^M\Bigl[g'(t)\mu(t)-\tfrac12 g'(t)^2 P(t)\Bigr]\dt
  \le\int_0^M\max_{s\ge0}\Bigl[s\,\mu(t)-\tfrac12 s^2 P(t)\Bigr]\dt
  =\int_0^M\frac{\mu(t)^2}{2P(t)}\dt,
\]
the pointwise maximum $\max_s(s\mu-\tfrac12 s^2P)=\mu^2/(2P)$ attained at
$s=\mu(t)/P(t)$. Equality holds iff $g'(t)=\mu(t)/P(t)$ a.e., which by
\eqref{eq:refcond} is an $L^\infty$ function, so $g_0\in\mathcal L$.
Multiplying by $2$ gives \eqref{eq:tmod-formula}.
\end{proof}

The next proposition is the isoperimetric statement that makes the
rearrangement decrease the Rayleigh quotient. It is the only place where
the Saint--Venant inequality enters.

\begin{proposition}[Isoperimetric inequality for $\Tmod$]\label{prop:isop-tmod}
Let $u\in H^1_0(\Om)\cap L^\infty$, $u\ge0$, satisfy \eqref{eq:refcond}.
If $B\subset\R^n$ is a ball with $T(B)=\Tmod(\Om;u)$, then $\abs B\le\abs\Om$,
with equality iff $\Om$ is a ball and $u$ is (a translate of) a radial
function.
\end{proposition}

\begin{proof}
$T(B)=\Tmod(\Om;u)\le T(\Om)$ by the remark after
Definition~\ref{def:tmod}. The Saint--Venant inequality \eqref{eq:sv-scale}
gives $T(\Om)\le T_n(\abs\Om/\omega_n)^{(n+2)/n}$, while
$T(B)=T_n(\abs B/\omega_n)^{(n+2)/n}$ by \eqref{eq:torsionball}. Hence
$(\abs B/\omega_n)^{(n+2)/n}\le(\abs\Om/\omega_n)^{(n+2)/n}$, i.e.\
$\abs B\le\abs\Om$. Equality forces $\Tmod(\Om;u)=T(\Om)$ and equality in
Saint--Venant, hence $\Om$ a ball; the maximiser $g_0\circ u$ of
\eqref{eq:tmod-def} then coincides with the (radial) torsion function, so
$u$ is radial. (\cite[Prop.~3.8]{Brasco2014}.)
\end{proof}

\subsubsection*{The Kohler--Jobin rearrangement}

We now build $u^*$. For $t\in[0,M]$ set $\Om_t=\{u>t\}$ and let
\begin{equation}\label{eq:Tt}
  \mathcal T(t):=\Tmod\bigl(\Om_t;(u-t)_+\bigr).
\end{equation}
Since $\abs{\{(u-t)_+>s\}}=\mu(t+s)$ and the coarea factor of $(u-t)_+$ at
level $s$ is $P(t+s)$, formula \eqref{eq:tmod-formula} gives the explicit
expression
\begin{equation}\label{eq:Tt-formula}
  \mathcal T(t)=\int_t^M\frac{\mu(\tau)^2}{P(\tau)}\,d\tau,
\end{equation}
so that $\mathcal T\in\mathrm{Lip}_{\rm loc}([0,M))$, $\mathcal T(M)=0$,
$\mathcal T(0)=\Tmod(\Om;u)=:T_\star$, and for a.e.\ $t$
\begin{equation}\label{eq:Tprime}
  \mathcal T'(t)=-\frac{\mu(t)^2}{P(t)}<0 .
\end{equation}
Thus $\mathcal T$ is a strictly decreasing bijection $[0,M]\to[0,T_\star]$.

\paragraph{Definition of $u^*$.}
Let $B=B_{R_\star}$ be the ball with $T(B)=T_\star$, i.e.\
$R_\star=(T_\star/T_n)^{1/(n+2)}$. For $\tau\in[0,T_\star]$ let $R(\tau)$ be
the radius with $T(B_{R(\tau)})=\tau$, i.e.\
\begin{equation}\label{eq:Rtau}
  R(\tau)=(\tau/T_n)^{1/(n+2)},\qquad
  \abs{B_{R(\tau)}}=\omega_n R(\tau)^n
  =\omega_n\,(\tau/T_n)^{n/(n+2)}.
\end{equation}
Let $\varphi:=\mathcal T^{-1}:[0,T_\star]\to[0,M]$ (so $\varphi(0)=M$,
$\varphi(T_\star)=0$, $\varphi'<0$). We define $u^*$ on $B$ to be the radial
decreasing function whose value on the sphere $\{\abs x=R(\tau)\}$ is
$\psi(\tau)$, where $\psi:[0,T_\star]\to[0,\infty)$ is the
locally Lipschitz (monotone, absolutely continuous) function determined by
\begin{equation}\label{eq:psi-ode}
  \bigl(-\psi'(\tau)\bigr)\,\mu^*\!\bigl(\psi(\tau)\bigr)
  =\bigl(-\varphi'(\tau)\bigr)\,\mu\!\bigl(\varphi(\tau)\bigr),
  \qquad \psi(T_\star)=0,
\end{equation}
$\mu^*$ denoting the distribution function of $u^*$. Concretely the
superlevel set $\{u^*>\psi(\tau)\}=B_{R(\tau)}$ has, by construction,
torsional rigidity exactly $\tau$, the value of the \emph{modified}
torsion that the corresponding superlevel set $\Om_{\varphi(\tau)}$ of $u$
carries. In particular $\{u^*>0\}=B=B_{R_\star}$, so
\begin{equation}\label{eq:TB}
  T(B)=T_\star=\Tmod(\Om;u)\ \ \le\ \ T(\Om).
\end{equation}
Since $\mu^*(\psi(\tau))=\abs{B_{R(\tau)}}=\omega_n R(\tau)^n$, equation
\eqref{eq:psi-ode} integrates to an explicit $\psi$; we shall only need the
two qualitative consequences \ref{prop:A}, \ref{prop:B}, proved next.

\begin{lemma}[Dirichlet energy is preserved]\label{lem:A}
$\displaystyle\int_B\abs{\nabla u^*}^2\dx=\int_\Om\abs{\nabla u}^2\dx$.
\end{lemma}

\begin{proof}
\emph{The original side.} By coarea and the change of variable
$t=\varphi(\tau)$ (Lemma~2.5 of~\cite{Brasco2014}, the one-dimensional area
formula for $\varphi=\mathcal T^{-1}\in\mathrm{Lip}_{\rm loc}$ with
$\varphi'<0$),
\begin{equation}\label{eq:dir-orig}
  \int_\Om\abs{\nabla u}^2\dx
  =\int_0^M P(t)\dt
  =\int_0^{T_\star}P(\varphi(\tau))\,(-\varphi'(\tau))\,d\tau .
\end{equation}
Differentiating $\mathcal T(\varphi(\tau))=\tau$ and using
\eqref{eq:Tprime} gives the key relation
\begin{equation}\label{eq:phiprime}
  -\varphi'(\tau)=\frac{1}{-\mathcal T'(\varphi(\tau))}
  =\frac{P(\varphi(\tau))}{\mu(\varphi(\tau))^2}.
\end{equation}
Hence, writing $P(\varphi)(-\varphi')=\mu(\varphi)^2\cdot
\frac{P(\varphi)}{\mu(\varphi)^2}\cdot(-\varphi')$ and using
\eqref{eq:phiprime} for the middle factor,
$P(\varphi)(-\varphi')=\mu(\varphi)^2(-\varphi')^2$, so \eqref{eq:dir-orig}
becomes
\begin{equation}\label{eq:dir-orig2}
  \int_\Om\abs{\nabla u}^2\dx
  =\int_0^{T_\star}\mu(\varphi(\tau))^2\,(-\varphi'(\tau))^2\,d\tau.
\end{equation}

\emph{The rearranged side.} For the radial $u^*$, the level set
$\{u^*=\psi(\tau)\}$ is the sphere of radius $R(\tau)$, on which
$\abs{\nabla u^*}$ is constant. A direct computation (differentiate
$u^*(x)=\psi(\tau)$ on $\abs x=R(\tau)$ and use $R(\tau)=(\tau/T_n)^{1/(n+2)}$)
gives, with $P^*(\tau'):=\int_{\{u^*=\tau'\}}\abs{\nabla u^*}\dH$ the coarea
factor of $u^*$,
\begin{equation}\label{eq:dir-rear}
  \int_B\abs{\nabla u^*}^2\dx
  =\int_0^{T_\star}\mu^*(\psi(\tau))^2\,(-\psi'(\tau))^2\,d\tau,
\end{equation}
the exact radial analogue of \eqref{eq:dir-orig2}; this is the
specialisation to $p=2$ of the identity~(4.4)
of~\cite{Brasco2014}, obtained there by the same coarea computation. Now
the defining relation \eqref{eq:psi-ode},
$\mu^*(\psi)(-\psi')=\mu(\varphi)(-\varphi')$, makes the integrands of
\eqref{eq:dir-rear} and \eqref{eq:dir-orig2} \emph{equal}:
\[
  \mu^*(\psi(\tau))^2(-\psi'(\tau))^2
  =\bigl[\mu^*(\psi(\tau))(-\psi'(\tau))\bigr]^2
  =\bigl[\mu(\varphi(\tau))(-\varphi'(\tau))\bigr]^2
  =\mu(\varphi(\tau))^2(-\varphi'(\tau))^2 .
\]
Hence $\int_B\abs{\nabla u^*}^2=\int_\Om\abs{\nabla u}^2$.
\end{proof}

\begin{lemma}[$L^2$ norm does not decrease]\label{lem:B}
$\displaystyle\int_B(u^*)^2\dx\ \ge\ \int_\Om u^2\dx$, and more generally
$\int_B f(u^*)\dx\ge\int_\Om f(u)\dx$ for every convex
$f:[0,\infty)\to[0,\infty)$ with $f(0)=0$, strictly if $f$ is strictly
convex unless $\Om$ is a ball and $u$ radial.
\end{lemma}

\begin{proof}
The point is the comparison of \emph{distribution functions} along the
rearrangement. Apply Proposition~\ref{prop:isop-tmod} not to $u$ but to the
truncations: for each $\tau\in[0,T_\star]$ the ball $B_{R(\tau)}$ has
$T(B_{R(\tau)})=\tau=\Tmod(\Om_{\varphi(\tau)};(u-\varphi(\tau))_+)$ by
construction, so Proposition~\ref{prop:isop-tmod} (applied to
$\Om_{\varphi(\tau)}$ with reference function $(u-\varphi(\tau))_+$) yields
\begin{equation}\label{eq:vol-cmp}
  \mu^*(\psi(\tau))=\abs{B_{R(\tau)}}\ \le\ \abs{\Om_{\varphi(\tau)}}
  =\mu(\varphi(\tau)),\qquad\tau\in[0,T_\star].
\end{equation}
Combining \eqref{eq:vol-cmp} with the defining relation
\eqref{eq:psi-ode}, $\mu^*(\psi)(-\psi')=\mu(\varphi)(-\varphi')$, gives
\begin{equation}\label{eq:psi-monotone}
  -\psi'(\tau)=\frac{\mu(\varphi(\tau))}{\mu^*(\psi(\tau))}\,(-\varphi'(\tau))
  \ \ge\ -\varphi'(\tau)\qquad\text{for a.e.\ }\tau,
\end{equation}
since $\mu(\varphi)\ge\mu^*(\psi)$. Integrating from $\tau$ to $T_\star$ and
using $\psi(T_\star)=\varphi(T_\star)=0$,
\begin{equation}\label{eq:psi-ge-phi}
  \psi(\tau)=\int_\tau^{T_\star}(-\psi'(s))\,ds
  \ \ge\ \int_\tau^{T_\star}(-\varphi'(s))\,ds=\varphi(\tau),
  \qquad\tau\in[0,T_\star].
\end{equation}
Now compute, for convex $f$ with $f(0)=0$ (so $f(u)=\int_0^u f'(t)\dt$ and
Cavalieri applies), using the change of variable $t=\varphi(\tau)$:
\begin{align*}
  \int_\Om f(u)\dx
  &=\int_0^M f'(t)\,\mu(t)\dt
  =\int_0^{T_\star}f'(\varphi(\tau))\,(-\varphi'(\tau))\,\mu(\varphi(\tau))\,d\tau\\
  &=\int_0^{T_\star}f'(\varphi(\tau))\,(-\psi'(\tau))\,\mu^*(\psi(\tau))\,d\tau
  \qquad\text{by \eqref{eq:psi-ode}}\\
  &\le\int_0^{T_\star}f'(\psi(\tau))\,(-\psi'(\tau))\,\mu^*(\psi(\tau))\,d\tau
  \qquad\text{by \eqref{eq:psi-ge-phi} and $f'$ nondecreasing}\\
  &=\int_0^M f'(s)\,\mu^*(s)\,ds=\int_B f(u^*)\dx,
\end{align*}
the last change of variable being $s=\psi(\tau)$. Taking $f(s)=s^2$ proves
$\int_B(u^*)^2\ge\int_\Om u^2$. If $f$ is strictly convex, equality in the
middle step forces $\psi=\varphi$ a.e., hence equality in
\eqref{eq:vol-cmp}, i.e.\ $\abs{B_{R(\tau)}}=\abs{\Om_{\varphi(\tau)}}$ for
a.e.\ $\tau$; by the equality case of the Saint--Venant inequality this
makes a.e.\ superlevel set $\Om_t=\{u>t\}$ a ball. The sets $\Om_t$ are
nested, so for $s>t$ the inclusion $\Om_s\subset\Om_t$ forces
$\abs{c_s-c_t}\le r_t-r_s$ on the centres and radii; as $t\downarrow0$ the
radii increase to a limit $r_\star$ and the centres are therefore Cauchy,
converging to some $c_\star$. Since the $\Om_t$ increase to $\{u>0\}=\Om$ (up
to a null set), $\Om=B(c_\star,r_\star)$; hence $\Om$ is a ball and $u$ is
radial, consistently with the equality case of
Proposition~\ref{prop:isop-tmod}.
\end{proof}

\begin{proof}[Proof of Theorem~\ref{thm:kj}]
Apply the construction to the first eigenfunction $u$ of $\Om$ (a valid
reference function by \eqref{eq:flux} and \ref{R:reg}). By
Lemmas~\ref{lem:A} and~\ref{lem:B} the rearranged $u^*\in H^1_0(B)$
satisfies $\int_B\abs{\nabla u^*}^2=\int_\Om\abs{\nabla u}^2$ and
$\int_B(u^*)^2\ge\int_\Om u^2$. By the Rayleigh principle on $B$
\eqref{eq:rayleigh} and these two facts,
\[
  \lambda_1(B)\le
  \frac{\int_B\abs{\nabla u^*}^2}{\int_B(u^*)^2}
  \le\frac{\int_\Om\abs{\nabla u}^2}{\int_\Om u^2}
  =\lambda_1(\Om).
\]
Multiplying $\lambda_1(B)\le\lambda_1(\Om)$ by $T(B)^\beta>0$ and using
$T(B)\le T(\Om)$ from \eqref{eq:TB} and the scale invariance
$\Psi(B)=\Psi(B_1)$:
\[
  T(B_1)L_n^\beta=\Psi(B)=T(B)\lambda_1(B)^\beta
  \le T(\Om)\lambda_1(\Om)^\beta=\Psi(\Om).
\]
This is \eqref{eq:kj} (for an arbitrary ball, $\Psi(B)=\Psi(B_1)$ on both
sides). Equality forces equality throughout, in particular in
Lemma~\ref{lem:B} with the strictly convex $f(s)=s^2$, hence $\Om$ is a
ball.
\end{proof}

\begin{remark}[On the sign that distinguishes Kohler--Jobin from Faber--Krahn]
\label{rem:sign}
The torsion-preserving rearrangement makes the ball
$B_{R(\tau)}$ \emph{smaller} than the superlevel set
$\Om_{\varphi(\tau)}$, \eqref{eq:vol-cmp} — the opposite of Schwarz
symmetrisation. It is precisely this that lets the Dirichlet energy be
\emph{exactly preserved} (Lemma~\ref{lem:A}) rather than merely decreased,
while the level-set volumes \emph{grow under the inverse rearrangement},
forcing $\int(u^*)^2\ge\int u^2$ (Lemma~\ref{lem:B}). The naive attempt to
compare $L^2$ masses level by level fails because the rearranged superlevel
balls are smaller; the correct comparison is through the
\emph{distribution-function inequality} \eqref{eq:psi-ge-phi}, driven by the
isoperimetric Proposition~\ref{prop:isop-tmod}.
\end{remark}

%---------------------------------------------------------------------
\subsection{Part 2: transfer to sharp Faber--Krahn stability}
\label{ssec:transfer}

We now use Theorem~\ref{thm:kj} as a tool. Throughout, $L:=\lambda_1(B^*)$,
$T_0:=T(B^*)$, $\beta=(n+2)/2$, $p:=1/\beta=2/(n+2)\in(0,1]$.

\begin{lemma}[Pointwise Kohler--Jobin transfer]\label{lem:ptwise}
For every open $\Om\subset\R^n$ with $\abs\Om=\abs{B^*}$,
\begin{equation}\label{eq:ptwise}
  \lambda_1(\Om)\ \ge\ L\,\Bigl(\frac{T_0}{T(\Om)}\Bigr)^{1/\beta}
  =L\,\Bigl(\frac{T_0}{T(\Om)}\Bigr)^{2/(n+2)} .
\end{equation}
\end{lemma}

\begin{proof}
Theorem~\ref{thm:kj} with $B=B^*$ gives
$T(\Om)\lambda_1(\Om)^\beta\ge T_0 L^\beta$; solve for $\lambda_1(\Om)$.
\end{proof}

By Saint--Venant \eqref{eq:sv-scale} at fixed volume, $T(\Om)\le T_0$, so
the torsion deficit $\sigma:=\sigma(\Om):=T_0-T(\Om)\ge0$. Put
$s:=\sigma/T_0\in[0,1)$; then \eqref{eq:ptwise} reads
\begin{equation}\label{eq:ptwise-s}
  \lambda_1(\Om)-L\ \ge\ L\bigl[(1-s)^{-p}-1\bigr],\qquad p=\tfrac{2}{n+2}.
\end{equation}

\begin{lemma}[An elementary one-variable inequality]\label{lem:onevar}
For $p\in(0,1]$ and $s\in[0,1)$, $\ (1-s)^{-p}-1\ge p\,s$.
\end{lemma}

\begin{proof}
$f(s)=(1-s)^{-p}-1$ has $f(0)=0$ and $f'(s)=p(1-s)^{-(p+1)}\ge p$ on
$[0,1)$ because $(1-s)^{-(p+1)}\ge1$. By the fundamental theorem of
calculus $f(s)=\int_0^s f'\ge ps$.
\end{proof}

\begin{lemma}[Small-deficit linearisation]\label{lem:small}
If $\abs\Om=\abs{B^*}$ and $\sigma\le T_0/2$, then
$\displaystyle\lambda_1(\Om)-L\ge\frac{2L}{(n+2)T_0}\,\sigma$.
\end{lemma}

\begin{proof}
With $s=\sigma/T_0\le1/2$, $p=2/(n+2)$, combine \eqref{eq:ptwise-s} and
Lemma~\ref{lem:onevar}:
$\lambda_1(\Om)-L\ge Lps=Lp\sigma/T_0=\tfrac{2L}{(n+2)T_0}\sigma$.
\end{proof}

\begin{lemma}[Large-deficit branch]\label{lem:large}
If $\abs\Om=\abs{B^*}$ and $\sigma>T_0/2$, then
$\lambda_1(\Om)-L\ge L(2^{2/(n+2)}-1)$, hence, since $\Fr(\Om)<2$,
\begin{equation}\label{eq:large}
  \lambda_1(\Om)-L\ \ge\ \frac{L(2^{2/(n+2)}-1)}{4}\,\Fr(\Om)^2 .
\end{equation}
\end{lemma}

\begin{proof}
$\sigma>T_0/2$ gives $T(\Om)<T_0/2$, so $T_0/T(\Om)>2$ and
\eqref{eq:ptwise} yields $\lambda_1(\Om)\ge L\,2^{2/(n+2)}$. Since
$\Fr(\Om)^2\le4$ at fixed volume, dividing by $4$ gives \eqref{eq:large}.
\end{proof}

\begin{theorem}[Sharp Faber--Krahn via Kohler--Jobin]\label{thm:fk}
By the global sharp Saint--Venant stability estimate of
Theorem~\ref{thm:SV-global}, with $c_{\mathrm{SV}}^{\mathrm{glob}}(n)>0$,
\begin{equation}\label{eq:sv-input}
  T(B^*)-T(\Om)=2\bigl(E(\Om)-E(B^*)\bigr)
  \ \ge\ 2\,c_{\mathrm{SV}}^{\mathrm{glob}}(n)
  \Bigl(\frac{\abs\Om}{\omega_n}\Bigr)^{(n+2)/n}\Fr(\Om)^2
\end{equation}
for all open bounded $\Om\subset\R^n$, $0<\abs\Om<\infty$. Then
\begin{equation}\label{eq:fk-out}
  \delta_{\mathrm{FK}}(\Om)
  =\Bigl(\frac{\abs\Om}{\omega_n}\Bigr)^{2/n}
   \bigl(\lambda_1(\Om)-\lambda_1(B^*)\bigr)
  \ \ge\ c_{\mathrm{FK}}(n)\,\Fr(\Om)^2,
\end{equation}
with the explicit constant
\begin{equation}\label{eq:cfk}
  c_{\mathrm{FK}}(n)=\min\!\left\{
  \frac{4\,L_n\,c_{\mathrm{SV}}^{\mathrm{glob}}(n)}{(n+2)\,T_n},\ \
  \frac{L_n\bigl(2^{2/(n+2)}-1\bigr)}{4}\right\},
\end{equation}
equivalently $\Fr(\Om)^2\le C_{\mathrm{FK}}(n)\,\delta_{\mathrm{FK}}(\Om)$
with $C_{\mathrm{FK}}(n)=1/c_{\mathrm{FK}}(n)$.
\end{theorem}

\begin{proof}
Both $\delta_{\mathrm{FK}}$ and $\Fr$ are scale invariant, and so is
\eqref{eq:sv-input}; rescale so that $\abs\Om=\omega_n$, whence $B^*=B_1$,
$L=L_n$, $T_0=T_n$, and the volume prefactor $(\abs\Om/\omega_n)^{2/n}=1$.
Split on the torsion deficit $\sigma=T_n-T(\Om)$.

\emph{Case 1 ($\sigma\le T_n/2$).} By Lemma~\ref{lem:small} and
\eqref{eq:sv-input} (which now reads
$\sigma\ge 2c_{\mathrm{SV}}^{\mathrm{glob}}(n)\Fr(\Om)^2$),
\[
  \lambda_1(\Om)-L_n
  \ge\frac{2L_n}{(n+2)T_n}\sigma
  \ge\frac{2L_n}{(n+2)T_n}\cdot 2c_{\mathrm{SV}}^{\mathrm{glob}}(n)\,\Fr(\Om)^2
  =\frac{4L_n c_{\mathrm{SV}}^{\mathrm{glob}}(n)}{(n+2)T_n}\,\Fr(\Om)^2 .
\]

\emph{Case 2 ($\sigma>T_n/2$).} By Lemma~\ref{lem:large},
$\lambda_1(\Om)-L_n\ge\frac{L_n(2^{2/(n+2)}-1)}{4}\Fr(\Om)^2$.

In both cases $\delta_{\mathrm{FK}}(\Om)=\lambda_1(\Om)-L_n\ge
c_{\mathrm{FK}}(n)\Fr(\Om)^2$ with $c_{\mathrm{FK}}(n)$ the minimum
\eqref{eq:cfk}. Undoing the scaling preserves the inequality since every
term is scale invariant.
\end{proof}

\begin{remark}[Sharpness of the exponent]\label{rem:sharp}
The exponent $2$ on $\Fr$ in \eqref{eq:fk-out} is optimal: smooth
ellipsoidal perturbations of the ball have
$\lambda_1(\Om)-L_n\asymp\Fr(\Om)^2$
(\cite{BraDeP2017,Brasco2014} and references therein). The transfer is
non-perturbative: no compactness, contradiction, or Taylor expansion is
used past the elementary Lemma~\ref{lem:onevar}.
\end{remark}

\begin{proof}[Proof of Theorem~\ref{thm:FK-main}]
For an open set of finite measure with $\lambda_1(\Om)=+\infty$ the inequality
is trivial, since its left side is then $+\infty$ while
$\Fr(\Om)^2\le4$.  Otherwise $\lambda_1(\Om)<\infty$, and
Theorem~\ref{thm:fk} applies: its Saint--Venant hypothesis
\eqref{eq:sv-input} is exactly the global estimate \eqref{eq:T-global} of
Theorem~\ref{thm:SV-global}, so \eqref{eq:fk-out} holds with
$c_{\mathrm{FK}}(n)$ as in \eqref{eq:cfk}.  This is the assertion of
Theorem~\ref{thm:FK-main}.
\end{proof}

%=====================================================================
